\subsection{The two center branches and their real period domains}

The center identities used below are identities of actual return maps
on specified real domains. In original coefficients the two branches
are
\[
 \mathcal C_r=\{a=c=d=0\},\qquad
 \mathcal C_c=\{m=-B,\ a=-c,\ d=c\}.
\]
For $\rho>0$ these become, respectively,
\[
 (\bar a,\bar c,\nu)=(0,0,0),\qquad
 (\bar m,\bar a,\nu)=(-\bar B,-\bar c,0).
\]
Their union is exactly the algebraic set
$\nu=\tau_1=\sigma_1=0$, because
$\tau_1=\bar a+\bar c$ and
$\sigma_1=\bar c(\bar B+\bar m)$.
We construct their real period domains by explicit first integrals with
compact level sets, fixing $\epsilon>0$ throughout this subsection.
These domains supply the actual center returns used in division.
The transition estimates of Section~\ref{rw:transitions} supply the
uniform finite jets as a gap or the scale vanishes.

\begin{proposition}[Real separated center]\label{rw:real-centers}
On $\mathcal C_c$, assume
\[
 b=1-\epsilon\bar B>0,\qquad a_0=1-\epsilon\bar c^2>0.
\]
Let $I_C$ be the component containing zero of
$\{\chi:A_C(\chi)>0\}$, where
\begin{equation}\label{rw:centers:separated-polynomial}
 A_C(\chi)=a_0+\bar c(1-2\epsilon\bar B)\chi
                       +b\bar B\chi^2,
\end{equation}
and let
\[
 J_C=\begin{cases}
 (-\epsilon,\infty),&\bar B\ge0,\\
 (-\epsilon,-1/\bar B),&\bar B<0.
 \end{cases}
\]
The change of state coordinates
\begin{equation}\label{rw:centers:separated-coordinate}
 \chi=-\frac{X+\bar cY}{1+\bar BY},\qquad
 X=-\bar cY-\chi(1+\bar BY)
\end{equation}
is a diffeomorphism on $1+\bar BY>0$. The inverse image
$\Omega_C$ of $I_C\times J_C$ is filled by the origin and periodic
ovals of the original field. More precisely,
\begin{equation}\label{rw:centers:separated-energy}
 \mathcal E_C(\chi,Y)
 =\int_0^\chi\frac{s}{A_C(s)}\,\dd s
   +\int_0^Y\frac{y}{(y+\epsilon)(1+\bar By)}\,\dd y
\end{equation}
is a real analytic first integral on this domain. It has one critical
point, a nondegenerate minimum at the origin, and compact sublevel
sets. Every level $\mathcal E_C=h>0$ is one smooth connected oval,
and the vector field is nonzero on it.
The construction includes $\bar B=0$ and both signs of the
quadratic discriminant. On the positive axis $Y=0$, membership in
this domain is precisely $-p\in I_C$.
\end{proposition}

\begin{proof}
The original scaled equations on this branch are exactly
\[
 X_\xi=-Y(1+\bar BY)+\epsilon(\bar BX^2-\bar cX),
 \qquad Y_\xi=(Y+\epsilon)(X+\bar cY).
\]
The change \eqref{rw:centers:separated-coordinate} is regular since
$1+\bar BY>0$; its $X$ derivative is
$-(1+\bar BY)^{-1}\ne0$. Differentiating its quotient and substituting
$X=-\bar cY-\chi(1+\bar BY)$ gives, without a time change,
\begin{equation}\label{rw:centers:separated-field}
 \chi_\xi=Y A_C(\chi),\qquad
 Y_\xi=-(Y+\epsilon)(1+\bar BY)\chi.
\end{equation}
The derivatives of the two potentials in
\eqref{rw:centers:separated-energy} are
$\chi/A_C(\chi)$ and $Y/((Y+\epsilon)(1+\bar BY))$.
Their Lie derivatives along \eqref{rw:centers:separated-field}
are $\chi Y$ and $-\chi Y$. Thus the first integral is exact.

Each potential decreases strictly towards zero from its left endpoint
and increases strictly from zero towards its right endpoint; its
minimum is zero. We verify divergence at all endpoints.
For the $Y$ potential the lower endpoint is $-\epsilon$ and
$1+\bar BY\to b>0$. Its derivative is asymptotic to
$-\epsilon/[b(Y+\epsilon)]$, so the potential tends to infinity.
If $\bar B<0$, the upper endpoint $-1/\bar B>0$ is a simple zero
of the other denominator factor and gives logarithmic divergence.
If $\bar B>0$, the potential at infinity is
$\bar B^{-1}\log Y+O(1)$; if $\bar B=0$, it is
$Y-\epsilon\log((Y+\epsilon)/\epsilon)$ and again diverges.

A finite endpoint of $I_C$ is a simple or double zero of $A_C$
and is not zero, because $A_C(0)=a_0>0$.
On the selected side its integrand has the sign of that endpoint
and absolute size bounded below by a positive constant times the
inverse distance to it. Integration from zero, in the appropriate
direction, therefore tends to positive infinity. At an infinite
endpoint with $\bar B>0$, the integrand is
$1/(b\bar B\chi)+O(\chi^{-2})$, giving logarithmic divergence.
With $\bar B=0$ and $\bar c\ne0$, its infinite-end asymptotic is
$1/\bar c$, giving linear growth of the potential in the outward
direction. With $\bar B=\bar c=0$ the potential is $\chi^2/2$.
Finally, if $\bar B<0$, the concave quadratic $A_C$ has two finite
roots of opposite signs, already covered above. These cases exhaust
the endpoints. Since both potentials are nonnegative and diverge at
their endpoints, every sublevel is compact in $I_C\times J_C$.

The only critical point is $(\chi,Y)=(0,0)$. In $(X,Y)$ its Hessian
is
\begin{equation}\label{rw:centers:separated-hessian}
 D^2\mathcal E_C(0,0)=\frac1{1-\epsilon\bar c^2}
       \begin{pmatrix}1&\bar c\\\bar c&1/\epsilon\end{pmatrix}>0.
\end{equation}
Write the potentials as $U(\chi),V(Y)$. For $h>0$, the two roots
$\chi_-<0<\chi_+$ of $U=h$ bound two graphs
$V(Y)=h-U(\chi)$, one on either side of $Y=0$.
They join at $(\chi_\pm,0)$ with nonzero gradient, forming one
smooth embedded oval. The field has its only equilibrium at the
origin, so its nonzero tangent field on each compact oval has a
finite positive period. At $Y=0$, $\chi=-X=-p$, giving the axis criterion.
\end{proof}

\begin{proposition}[The reversible well]\label{rw:centers:reversible}
On $\mathcal C_r$, assume $1+\epsilon\bar m>0$ and
$|\epsilon\bar B|<1/4$. Put $\theta=\epsilon\bar B$,
$Z=Y+\epsilon>0$ and
\begin{equation}\label{rw:centers:reversible-energy}
 \mathcal E_r(X,Y)=X^2 Z^{-2\theta}+V_r(Y),\qquad
 V_r(Y)=2\int_0^Y u(1-\bar m u)(u+\epsilon)^{-1-2\theta}\,\dd u.
\end{equation}
Let $Y_s=1/\bar m$ if $\bar m>0$, and $Y_s=\infty$ otherwise.
On $J_r=(-\epsilon,Y_s)$ define
\[
 h_- =\lim_{Y\downarrow-\epsilon}V_r(Y),\qquad
 h_+ =\lim_{Y\uparrow Y_s}V_r(Y),\qquad h_* =\min(h_-,h_+).
\]
Both limits are positive or infinite. The set
$\{Y\in J_r,\ 0<\mathcal E_r<h_*\}$ is a period annulus:
every level in this range is one compact periodic oval. If
$\bar B\ge0$ and $\bar m\le0$, then $h_*=\infty$ and this
annulus fills the half-plane $Y> -\epsilon$, apart from its center.
If $\bar B=0$ and $\bar m>0$, its upper boundary is the
homoclinic loop through $(0,1/\bar m)$.
\end{proposition}

\begin{proof}
Here the field is
$X_\xi=Y(\bar mY-1)+\theta X^2$,
$Y_\xi=(Y+\epsilon)X$.
Differentiating $X^2Z^{-2\theta}$ cancels its cubic terms and gives
$2X Y(\bar mY-1)Z^{-2\theta}$; the derivative of $V_r$
is its negative. This proves the first integral.
On $J_r$ the factor $1-\bar mY$ is positive. Thus $V_r$ decreases
strictly from its left endpoint to its minimum zero at $Y=0$ and
increases strictly thereafter. Its second derivative at zero is
$2\epsilon^{-1-2\theta}>0$.
At the lower boundary its derivative has a negative nonzero multiple
of $(Y+\epsilon)^{-1-2\theta}$ as its leading term. Hence
$h_- =\infty$ when $\theta\ge0$; when $\theta<0$ the finite
limit is positive. If $\bar m>0$, $h_+$ is finite and positive.
If $\bar m\le0$, the restriction $|\theta|<1/4$ makes the
positive-end integral divergent, both for the linear and quadratic
numerator.

For $0<h<h_*$, the regular coordinate $u=XZ^{-\theta}$ changes
the energy into $u^2+V_r(Y)$. Its level is one smooth compact oval,
by the same two-graph construction as above. The field vanishes only
at the origin in $J_r$, so the oval is periodic. The restriction
$h<h_*$ retains compactness when $\bar B<0$ or $\bar m>0$.

When $\bar B=0$ and $\bar m>0$, $h_- =\infty$ and the upper
critical point of $V_r$ at $Y_s$ is a nondegenerate maximum.
The energy $h_+=V_r(Y_s)$ has a negative turning point and two
arcs meeting at the saddle $(0,Y_s)$. The field has no other zero
on those arcs. They form the bounding homoclinic loop of the central
well. For $\bar B\ge0$, $\bar m\le0$, both limiting energies
are infinite, proving the full half-plane assertion.
\end{proof}

\subsection{The three-saddle triangle and the signed weak boundary}

The separated integral covers either sign of $\bar B$. We now extract
two boundary descriptions needed later: exact upper power coordinates
on the three-saddle triangle, and a common finite strong and upper weak
intersection through $\bar B=0$.

\begin{lemma}[The triangle and its upper saddle]\label{rw:centers:triangle}
On $\mathcal C_c$ write $\bar B=-\beta<0$,
$\bar m=\beta$, $\bar c=-\bar a$, and assume
$1-\epsilon\bar a^2>0$. Set
\[
 b=1+\epsilon\beta,\qquad
 \Delta=\sqrt{\bar a^2+4\beta b},\qquad
 k_\pm=\frac{\bar a\pm\Delta}{2b},
 \qquad
 p_+=\frac{\bar a-k_-}{\beta},\quad
 p_- =\frac{k_+-\bar a}{\beta}.
\]
Then $p_\pm>0$ and the period domain $\Omega_C$ is the triangle
\begin{equation}\label{rw:centers:triangle-domain}
 \mathcal T=\{L_0=Y+\epsilon>0,\quad
 L_+=p_+-X+k_-Y>0,\quad
 L_-=p_-+X-k_+Y>0\}.
\end{equation}
Its vertices are
$((\bar a\pm\Delta)/(2\beta),-\epsilon)$ and
$(\bar a/\beta,1/\beta)$.
The positive function
\begin{equation}\label{rw:centers:triangle-integral}
 \mathcal H=L_0L_+^{\lambda_+}L_-^{\lambda_-},\qquad
 \lambda_\pm=\frac{\Delta\pm\bar a(1+2\epsilon\beta)}
                         {2\epsilon\beta\Delta}
\end{equation}
is a first integral on $\mathcal T$.
Near the upper vertex the analytic coordinates
\[
 U=L_+L_0^{\epsilon\beta},\qquad
 V=L_-L_0^{\epsilon\beta}
\]
have a regular inverse and give the exact orbital equation
\begin{equation}\label{rw:centers:triangle-power}
 \frac{\dd V}{\dd U}=-r\frac VU,
 \qquad r=\frac{\lambda_+}{\lambda_-}=\frac{p_+}{p_-}>0.
\end{equation}
On compact $\beta$ intervals bounded away from zero, these local
coordinate changes have bounded finite jets for small scales and
bounded $\bar a$.
\end{lemma}

\begin{proof}
The identity
\[
 \Delta^2-\bar a^2(1+2\epsilon\beta)^2
       =4\beta b(1-\epsilon\bar a^2)>0
\]
proves positivity of $p_\pm$ and $\lambda_\pm$.
The lines in \eqref{rw:centers:triangle-domain} meet at the stated
vertices. In the separated coordinate the roots of $A_C$ are
$-p_+$ and $p_-$; its leading coefficient is negative.
For $Y<1/\beta$, the inequalities
$-p_+<\chi<p_-$ translate exactly into $L_+>0,L_->0$.
Thus the triangle agrees with $\Omega_C$.

Direct differentiation of its three affine factors gives their
cofactors
\[
 X-\bar aY,\qquad
 -\epsilon\beta X-bk_-Y,\qquad
 -\epsilon\beta X-bk_+Y.
\]
The identities
$\lambda_++\lambda_-=1/(\epsilon\beta)$ and
$\lambda_+k_-+\lambda_-k_+=-\bar a/b$
make their weighted sum zero. This proves the first integral.
Its logarithm is strictly concave and tends to negative infinity
at the boundary. At the origin
\[
 \frac{\lambda_+}{p_+}=\frac{\lambda_-}{p_-}
                  =\frac{b}{\epsilon\Delta},
\]
so both components of its gradient vanish. The energy
$-\log(\mathcal H/\mathcal H(0,0))$ is consequently proper and
strictly convex on the triangle. Its positive levels give the same
periodic ovals described by Proposition~\ref{rw:real-centers}.

At the upper vertex $L_0=\epsilon+1/\beta>0$ and the two other
factors have independent gradients. Multiplication by the positive
analytic unit $L_0^{\epsilon\beta}$ preserves their invertible
Jacobian. Since $\epsilon\beta(\lambda_++\lambda_-)=1$,
\[
 \mathcal H=U^{\lambda_+}V^{\lambda_-}.
\]
Differentiation along the nonzero saddle orbit gives
\eqref{rw:centers:triangle-power}. The ratio equals $p_+/p_-$ by
the preceding gradient identity. Near this finite positive-height
vertex, $L_0$ is bounded away from zero on each stated compact block;
all finite derivatives of the units and their inverses are bounded.
\end{proof}

\begin{lemma}[A common real weak boundary]\label{rw:centers:weak-boundary}
Fix $\bar c>c_*>0$ in a compact interval and let $\bar B$ vary
in a neighborhood of zero on $\mathcal C_c$. For sufficiently
small $\epsilon>0$ define
\begin{equation}\label{rw:centers:weak-intersection}
 \Delta_b=\sqrt{\bar c^2-4b\bar B},\qquad
 k_s=\frac{-\bar c-\Delta_b}{2b},\qquad
 \widehat p_C=\frac{2(1-\epsilon\bar c^2)}
                 {\bar c(1-2\epsilon\bar B)+\Delta_b}.
\end{equation}
These are regular through $\bar B=0$ and
$\bar B\widehat p_C=\bar c+k_s$.
The invariant line
$L=\widehat p_C-X+k_sY=0$ is the boundary of $\Omega_C$
met by its positive finite axis inputs. Thus inputs
$p=\widehat p_C-t$, $0<t\ll1$, belong to the period domain.
In the reversed upper chart $R=1/Y$, $y=X/Y-k_s$, the same line
is $y=\widehat p_CR$, and its restricted equation is
\begin{equation}\label{rw:centers:weak-restricted}
 \dot R=\widehat p_CR(R+\bar B)(1+\epsilon R),\qquad
 \gamma_C=\bar B\widehat p_C.
\end{equation}
In particular $\partial_{\bar B}\gamma_C>c_1>0$ on a common
smaller block. The finite strong and upper weak intersections on
this center are the same actual point $p=\widehat p_C$.
\end{lemma}

\begin{proof}
The discriminant of $A_C$ is
$\bar c^2-4b\bar B$. Rationalizing its root nearest zero on the
negative side gives $-\widehat p_C$. The displayed denominator
has a fixed positive margin. Substitution in
$bk_s^2+\bar c k_s+\bar B=0$ proves
$\bar B\widehat p_C=\bar c+k_s$, including its continuous
value at $\bar B=0$. The equation $\chi=-\widehat p_C$
from \eqref{rw:centers:separated-coordinate} is precisely $L=0$.
Its cofactor is $\epsilon\bar BX-bk_sY$, so the line is
invariant under the original field.

The first positive zero on $Y=-\epsilon$ of the horizontal field
is
\[
 p_\epsilon=\frac{2b}{\bar c+\Delta_b}.
\]
The identity $p_\epsilon-\epsilon\bar c=b\widehat p_C$
shows that this zero lies on $L=0$. Its branch up to $Y=0$
is the same finite strong boundary just identified by the scalar
center domain. In the reversed upper chart substitute
$y=\widehat p_CR$ in the exact projective field and use
$\bar c+k_s=\bar B\widehat p_C$. This gives
\eqref{rw:centers:weak-restricted} on the actual weak face. At $\bar B=0$,
\[
 \partial_{\bar B}\gamma_C=\widehat p_C
                =\frac{1-\epsilon\bar c^2}{\bar c}>0.
\]
Continuity on the compact block proves the uniform positive margin.
The actual invariant line can be preserved in any auxiliary extension
used to construct its weak graph, so the two constructions retain this
same intersection.
\end{proof}

For any oval in either of the stated domains, the two paths from its
axis intersections meet at the same original central height
$z=1+Y/\epsilon>0$. Their complete central maps therefore agree
on every common section interval parametrizing these ovals.
This identity holds throughout each stated real period domain. The
constructions below provide the signed extensions and coefficient
segments needed for division. The triangle gives an exact local saddle
power, while the reversible well retains its own period annulus; the
separated domain remains the one used across signed $\bar B$ faces.
